\section{Records, artifacts, and replay}\label{app:records}

\subsection{An interoperable record specification}

The following fields specify the information an RRE implementation should expose. The case-study files implement these concepts through a charter, typed graph records, pinned evidence, event records, and version deltas. The underlying atlas's extraction schema is described in \cref{app:complexity-graphs}.

\begin{longtable}{@{}p{2.0cm}p{4.0cm}p{7.8cm}@{}}
\toprule
Record & Core fields & Semantic requirement \\
\midrule
\endhead
Claim & ID, revision, statement, role, scope, source & Quantifiers and parameter regime belong to the statement; reconstructed text and source evidence remain distinguishable. \\
Inference & ID, premise IDs, conclusion, scope, justification & All premises are jointly required. Alternative proofs and conditional branches receive separate records. \\
Target & ID, parent, statement, objective link, bounds, status & Success must meet the stated quantitative requirement; partial results remain separately usable. \\
Objective & ID, origin, rationale, acceptance conditions, revision & Preserve the original objective and any prospectively recorded family changes. \\
Evidence & Type, location, hash, inputs, environment & Distinguish exact computations, source assertions, written proofs, and formal certificates. \\
Review & Reviewer role, coverage, method, findings & Report the actual scope of checking. \\
Event & ID, time, actor, target, action, inputs, outcome, change & Separate the actor of an action from the authorship of its imported sources. \\
Assessment & Result, significance, rationale, literature comparison & Assess resolved obligations separately from proof status. \\
Run & Mode, objective, initial graph, sources, resources, stop reason & Retain acceptance conditions and the actual stopping decision. \\
Strategy revision & Method versions, evaluation tasks, costs, decision & Evaluate changes to the research procedure separately from its mathematical outcomes. \\
\bottomrule
\end{longtable}

A node ID remains stable while a version identifies its exact content. A source anchor records a file or public source version, its location, and a hash where available. Graph evidence in this run pins local mathematical notes, certificates, checks, and reading records. Public references specify the consulted versions; a hash of a local reading note is not presented as a hash of the cited PDF.

For example, the exact two-term theorem has an exhaustive-case inference. Its inputs are the lower bound, the geometric dichotomy, and two conditional upper-bound routes. The boundary route depends on the projection to $W_3$, the base certificate, and balanced transfer. The separate $W_n$ inference uses the explicit family contraction and does not import the general geometric classification. Both traces retain the original-parameter degree convention. The historical event that proposed transfer is linked to these claims as motivation rather than inserted as a mathematical premise. The normalized inequality in \cref{cor:degree-transfer} and the matrix endpoint of \cref{thm:main} make further consequences of the recorded lemmas explicit. The six graph versions retain their original statements and scopes as the exploration history.

\subsection{Artifacts and portable replay}

The companion contains the six component graphs and their deltas, the charter, the source pins, 44 event records, the mathematical notes, exact certificates, and computational reports. A source manifest records the included bytes and their SHA-256 hashes. The portable replay program checks the saved rational certificates directly in the original variables and replays the transferred $W_n$ identities. It also validates graph identifiers, scopes, acyclic proof dependencies, version links, and event hashes. The original research utilities and saved reports are retained for comparison; replay writes its own report.

\begin{table}[htbp]
\centering\small
\begin{tabularx}{\linewidth}{@{}p{5.0cm}Y@{}}
\toprule
Artifact & What it makes reproducible or inspectable \\
\midrule
Charter and initial source pins & Prospective objective, degree convention, inherited inputs \\
Six graphs and version deltas & Incremental additions, revisions, and retained open questions \\
Event history and replay index & Targets, actions, outcomes, actor attribution, and artifact versions \\
Mathematical notes and certificate table & Uniform arguments and the complete degree-9 identity \\
Rational certificates for three single-block pencils & Raw-variable unit identities of degrees 9, 12, and 15 \\
Transfer and contraction checks & All-order construction tested on explicit finite inputs \\
Uniform commutator checks & Symbolic identities and singular-parameter controls \\
Prior-work comparison & Classical inputs and the specific certificate-degree contribution \\
Viewer, captures, and figure map & Connection between displayed claims and saved graph versions \\
Dedicated discovery exports & Prerequisite paths and evidence for each principal theorem \\
\bottomrule
\end{tabularx}
\caption{Principal artifacts of the certificate-degree case study. Exact filenames and replay commands are supplied in the companion documentation.}
\end{table}

The rational certificate files are replayed by expanding each normalized orbit monomial into original parameter monomials, multiplying by its original coefficient equation, and collecting coefficients over $\Q$. The result must be exactly the constant one, and the largest product degree is checked. This verification uses the certificate rather than repeating the elimination search. The appendix's 19-row table is also compared against the certificate file, so the printed mathematical object and its machine-readable counterpart agree.

The all-order transfer is proved in \cref{lem:transfer}; its finite replay uses the same base coefficients with the contracted parameter factors retained. Source-hash validation checks which recorded artifact version supports each claim. Mathematical checking and provenance checking produce separate outcomes. Replaying the saved artifacts does not require restarting the original research agent or regenerating its target choices.

\subsection{Event-level effects}\label{app:event-effects}

The table groups all 44 saved records by their effect. The raw display IDs are preserved. Two concurrent computation-start records both used E0022; their distinct event hashes identify them uniquely, and the replay index assigns unique sequential identities. No mathematical result depends on the duplicated display number.

\small
\begin{longtable}{@{}p{2.7cm}p{11.3cm}@{}}
\toprule
Recorded events & Effect on the investigation and graph \\
\midrule
\endhead
E0001--E0004 & Adopt the objective and initial targets, seal v001, and publish its atlas component. \\
E0005--E0008 & Perform the first degree tests, verify the matrix-pencil baseline, extend the lower functional to unequal dimensions, and prove complete symmetry reduction. \\
E0009--E0012 & Generate the all-size sharp-bound and target-grading questions, publish v002, and replay the rational certificates by raw expansion. \\
E0013--E0015 & Use target grading, begin the rational lift of the degree-15 test, and save its exact certificate and expansion. \\
E0016--E0019 & Generate the uniform commutator target, check its identities, record the all-$m$ theorem, and generate the remaining endpoint question. \\
E0020--E0021 & Seal and publish v003 with the uniform theorem and its proof obligations. \\
E0022 (two records), E0024--E0026 & Start two multiple-block screens, save their distinct outcomes, and preserve and reconcile the concurrent event-log fork. \\
E0027--E0031 & Publish v004, correct an export path-length issue, export the pencil theorem, and save the open exact-endpoint checkpoint. \\
E0032--E0033 & Clarify the acceptance policy without weakening the mathematical objective; prospectively generate the balanced-transfer and full CP-2 target. \\
E0034--E0036 & Replay transfer and contraction maps, then save the all-order exact theorem and its $W_n$ specialization. \\
E0037--E0038 & Seal and publish v005, including the classical imported supports and the two conditional routes. \\
E0039--E0042 & Audit the theorem's scope, degree accounting, and case coverage; assess the original criteria; seal and publish v006. \\
E0043--E0044 & Export the exact theorem's prerequisite graph and preserve the completed checkpoint and replay outline. \\
\bottomrule
\end{longtable}
\normalsize

The event fork preserves both original records and their common predecessor. A reconciliation record explains the fork, and subsequent writes use serialization. Playback follows the recorded file order with unique event hashes; the raw log remains unchanged. This is an explicit part of the provenance rather than a reconstructed single chain with an event silently removed.

The six graph publications are actual saved revisions. Each summarizes the claims admitted by that checkpoint; it is not an invented graph state for every intervening action. The final proof can therefore be read in dependency order, while the events preserve the distinct order in which targets, checks, and proofs were recorded.

\subsection{Visualization and evidence links}

The companion supplies the original graph data embedded in offline viewer pages, one per version. Selecting an inference exposes all its premises and its justification; selecting a claim exposes its statement, attached scopes, status, and evidence. Prerequisite traversal follows producing inferences and attached scopes. Navigation links are optional and do not participate in that traversal.

The printed overview uses all six saved graphs, while the inspector image shows a selected record and its evidence. The browser validation report records graph hashes and counts and checks the prerequisite closures of the two-term theorem and $W_n$. The capture routine also checks that displayed node text stays within its cards. The timeline is an explanatory diagram keyed to the same versions and events; its arrows describe the emergence of research targets rather than proof implications.

The versioned viewer extends the earlier graph-comparison paradigm by providing a separate view for each of the six saved revisions. The companion pairs these views with the corresponding graph data and evidence.
